\chapter{Results and Discussion}
\label{ch:results}

\section{Delivered Capabilities}

\noindent Table~\ref{tab:results} maps the objectives of
Chapter~\ref{ch:introduction} to the delivered, verified capabilities.

\begin{table}[h]
\centering
\small
\begin{tabular}{@{}p{2.9cm}p{11.4cm}@{}}
\toprule
\textbf{Objective} & \textbf{Delivered capability} \\
\midrule
Tenancy and identity & Multi-institute deployment with memberships, roles,
and a permission catalogue; every query tenant-scoped by the enforcement
layers. \\
Academic structure & Years, classes, divisions, subjects, chapters, topics,
offerings, teacher assignments, student placements/enrollments; academic scope
derived from current DB state. \\
Materials and OCR & Chunked upload; distributed OCR (embedded-text-first,
PaddleOCR fallback); page corrections; coordinator-owned chunk leasing. \\
Syllabus intelligence & AI analysis into per-subject structure and context
with PROPOSED$\rightarrow$CONFIRMED lifecycle and per-subject versioning. \\
Aided content creation & Typed, versioned content items generated from
eligible sources with provenance and in-place regeneration. \\
Question assets & Data-driven question types (11 global seeds), an approved
question bank, pattern-aware generation, paper patterns, and question papers. \\
Examination & Blueprint-driven assessments with schedule validation,
snapshot attempts, server-side objective evaluation, and aggregate analytics. \\
Practice & Ungraded flashcard and question sessions with instant per-item
feedback. \\
Export & Single-model PDF (Chromium), DOCX, XLSX rendering for content,
questions, papers, patterns, and results; preview equals export. \\
Security & Session-bound tokens, strict rotation with lineage revocation,
double-submit CSRF, DB-fresh authorisation, academic scope, and a
documentation-verified audit with no open high findings. \\
\bottomrule
\end{tabular}
\caption{Objectives versus delivered capabilities.}
\label{tab:results}
\end{table}

\section{Validation Summary}

\noindent The implementation is supported by 226 passing API unit tests,
eleven database-backed integration suites, strict type-checking and linting
across TypeScript and Python, a standing user-validation log driven by a
seeded demo dataset, and a security audit whose findings are all remediated
and covered by regression tests.

\section{Discussion}

\noindent Three architectural choices deserve discussion.

\subsection{Modular monolith over microservices}

Keeping the business modules in one deployable unit means authorization
evaluation, tenancy scoping, and database transactions stay simple and local.
The two genuinely independent workloads --- OCR and AI generation --- were
pushed out of the monolith to dedicated components behind queues, which is
where elasticity actually matters. This matches the project's standing
decision to avoid premature microservices.

\subsection{Permission + scope instead of pure RBAC}

Nesting academic scope under the permission model closes the classic RBAC
leak: a teacher with \texttt{questions.update} still cannot touch a question in
a subject they do not teach, because the scope gate evaluates the resource's
subject against the teacher's assignment-derived subject set. The cost is a
per-query scope predicate, which is cheap relative to the guarantee it gives.
The 404-for-reads / 403-for-writes asymmetry keeps resource existence hidden
from unauthorized actors.

\subsection{Coordinator-owned OCR versus queue-based work}

OCR workers are heterogenous and expected to be occasionally offline, so chunk
ownership is expressed as a \emph{lease} with heartbeat renewal and
sweep-based reclamation rather than an unbounded queue message. This design was
chosen over retried queue delivery and validated against a real database in
the OCR integration suite. A side effect is that the OCR pipeline can be
incrementally scaled by adding workers with no coordinator downtime.

\section{Limitations}

\noindent The following limitations are accepted and documented in the
repository:

\begin{itemize}
  \item \textbf{Selection randomness.} Pattern-to-paper selection is
  deterministic; randomization within a bucket is not implemented.
  \item \textbf{Question-paper immutability.} The question-paper junction
  keeps a live foreign key to the bank question; a later edit to that question
  can leak into paper exports (attempts, by contrast, are true snapshots).
  \item \textbf{Provenance reads.} The \texttt{source\_pattern\_id} column is
  write-only; there is no read path exposing a question's origin pattern in the
  API or web client.
  \item \textbf{Attempt-N-of-M.} Attempt numbering and max-attempt limits are
  presentation metadata only; all linked questions are always evaluated.
  \item \textbf{Duplication detection.} Within-batch exact-stem de-duplication
  and a prompt instruction are used; there is no similarity/embedding scan
  across the bank.
  \item \textbf{Scope granularity.} Academic scope is currently subject-set
  based; the designed offering-level scope (\texttt{offeringId} column) is
  not yet in the schema.
  \item \textbf{Rate limiting.} The global throttle is in-memory per instance;
  a distributed (Redis-backed) limiter is deferred until cross-instance
  rate limiting is needed.
  \item \textbf{OCR/AI quality metrics.} Extraction accuracy and generated
  content are validated structurally; quantitative quality metrics are not
  produced.
\end{itemize}

Further deferred scope --- online answer sheets (FORM), OMR, on-screen
marking, automated text grading, CRUD for a full \texttt{SUPER\_ADMIN}
institute lifecycle UI, and web-level end-to-end automation --- is documented
in Chapter~\ref{ch:conclusion}.